
This study designed and validated a statistical methodology for measuring ML dataset documentation compliance at initial release (T0+90 days) using a 3-tier temporal detection protocol. Applied to synthetic data matching expected distributions, the methodology successfully detected: (1) severe documentation gaps (7\% compliance, CI: [3.4\%, 13.8\%]), (2) non-uniform component deficiencies (licensing 27\% vs. data context 77\%), and (3) mechanism specificity (commit velocity $\rho$ = 0.951 vs. contributors $\rho$ = 0.028).

The validation demonstrates that the study design can establish temporal precedence for documentation gaps and identify specific correlates of documentation quality when applied to real repositories. The 3-tier T0 detection protocol enables retrospective temporal analysis without requiring prospective longitudinal data collection. The mechanism specificity tests successfully isolate commit velocity from contributor diversity and issue responsiveness.

\textbf{Critical caveat:} All numerical results are based on synthetic data designed to validate statistical methodology and gate logic. The actual compliance rates, correlation magnitudes, and component hierarchies require empirical confirmation through production deployment with HuggingFace Hub API sampling, GitHub API activity collection, and manual DCS\_3 coding.

\subsubsection{Future Work}

\textbf{FW1: Production Deployment with Real Data} --- Replace synthetic data with actual HuggingFace Hub API sampling and GitHub API activity collection to confirm or revise the 7\% compliance rate and $\rho$ = 0.951 correlation magnitude.

\textbf{FW2: Longitudinal Causal Test} --- Measure DCS\_3 at T0, T+30, T+60, T+90, T+180 for the same repositories to test whether commit spikes at time t precede DCS improvements from t to t+30, establishing temporal precedence for potential causation.

\textbf{FW3: Multi-Platform Replication} --- Apply DCS\_3 protocol to N=100 datasets each from Papers with Code, OpenML, and Zenodo to test whether patterns observed in synthetic data (if confirmed on HuggingFace) are platform-specific or field-wide.

\textbf{FW4: Intervention Testing} --- If real data confirms low licensing compliance, conduct randomized A/B test of automated licensing template prompts vs. standard upload to measure causal intervention effects.

The methodology validated in this study provides a foundation for empirical measurement of documentation practices in ML repositories. The next step is deployment on real data to determine whether the synthetic patterns reflect actual repository behavior.
